\chapter{A Concrete Naming Certificate for $\RegularCauchyMinModulusUp$}
\label{ch:concrete-instances-regular-cauchy-min-modulus-namecert}
\origin{human}

Following Bishop and Bridges constructive analysis, the $\RegularCauchyMinModulusUp$ packet records the finite scheduling surface for taking the pointwise minimum of two displayed regular-Cauchy modulus schedules. It sits between the shared-modulus meet packet of \autoref{ch:concrete-instances-cauchymodulusmeet-namecert}, the finite stream windows of \autoref{ch:concrete-instances-streamname-namecert}, the regular rational readback boundary of \autoref{ch:concrete-instances-regseqrat-namecert}, the dyadic tolerance ledger of \autoref{ch:concrete-instances-dyadicratcore-namecert}, and the terminal real-name seal of \autoref{ch:concrete-instances-real-namecert}. The packet is a finite L10 Real-and-Completion route: it names the displayed minimum-modulus schedule, not a selected quotient stream, an ambient completeness theorem, or a choice principle over all precision requests.

\begin{definition}[Regular Cauchy minimum-modulus carrier]\label{def:regular-cauchy-min-modulus-carrier}
A $\RegularCauchyMinModulusUp$ carrier is a finite $\BHist$ packet
$$
R=(S_0,S_1,M,m,W,Q,D,E,H,C,P,N)
$$
with the following displayed rows:
\begin{enumerate}
\item $S_0$ and $S_1$ are the two regular-Cauchy schedules whose finite precision requests are to be synchronized.
\item $M$ is the $\CauchyModulusMeetUp$ row that displays the common meet surface inherited from the two source schedules.
\item $m$ is the finite minimum row, recording the requestwise smaller admitted threshold after the two source schedules and the meet row have been read.
\item $W$ is the $\StreamNameUp$ window ledger naming exactly which source windows are inspected by the minimum schedule.
\item $Q$ is the $\RegSeqRatUp$ readback row attached to the same windows and the same displayed thresholds.
\item $D$ is the $\DyadicRatCoreUp$ tolerance ledger storing denominator budgets, dyadic radii, slack comparisons, and the finite inequalities used by $M$ and $m$.
\item $E$ is the terminal $\RealUp$ seal reached only after $S_0,S_1,M,m,W,Q,D$ have been displayed.
\item $H$ stores componentwise $\hsame$ transport for the two schedules, the meet row, the minimum row, the stream windows, the regular-rational readback, the dyadic ledger, the real seal, the provenance row, and the local name.
\item $C$ stores displayed $\Cont$ replay from a finite precision request through $S_0,S_1,M,m,W,Q,D,E,H$.
\item $P$ is the $\Pkg$ provenance over $\RegularCauchyMinModulusUp$, $\CauchyModulusMeetUp$, $\StreamNameUp$, $\RegSeqRatUp$, $\DyadicRatCoreUp$, $\RealUp$, $\BHist$, $\hsame$, $\Cont$, and $\NameCert$.
\item $N$ is the local $\NameCert_{\RegularCauchyMinModulusUp}$ row.
\end{enumerate}
The classifier compares accepted packets only through these rows. It has no coordinate for quotient stream equality, a selected limit, host real equality, countable choice, or ambient $\RealUp$ completeness.
\end{definition}

\begin{theorem}[Regular Cauchy minimum-modulus NameCert obligations]\label{thm:regular-cauchy-min-modulus-namecert-obligations}
For every accepted $\RegularCauchyMinModulusUp$ carrier
$$
R=(S_0,S_1,M,m,W,Q,D,E,H,C,P,N),
$$
the local $\NameCert_{\RegularCauchyMinModulusUp}$ surface consists exactly of the following finite obligation rows:
\begin{enumerate}
\item carrier admission for the two displayed regular-Cauchy schedules, the $\CauchyModulusMeetUp$ row, the finite minimum row, stream windows, regular rational readback, dyadic tolerance ledger, terminal real seal, transport, replay, provenance, and local naming rows;
\item meet compatibility, requiring $M$ to be read through the sibling shared-modulus surface of \autoref{ch:concrete-instances-cauchymodulusmeet-namecert} before the minimum row $m$ is consumed;
\item finite minimum exactness, requiring each request in $m$ to name the smaller displayed threshold among the two transported source schedules after the dyadic comparisons in $D$ have been read;
\item readback stability, requiring the same $\StreamNameUp$ windows $W$ and $\RegSeqRatUp$ row $Q$ to be replayed under componentwise $\hsame$ transport;
\item ledger rows for $\Cont$ replay, $\Pkg$ provenance, the $\RealUp$ seal handoff, and the local naming certificate.
\end{enumerate}
No accepted obligation row is stored outside $S_0,S_1,M,m,W,Q,D,E,H,C,P,N$.
\end{theorem}
\begin{proof}
Carrier admission is projection from \autoref{def:regular-cauchy-min-modulus-carrier}. The accepted packet displays the two schedules, the meet row, the minimum row, finite windows, regular readback, dyadic ledger, Real seal, transport, replay, provenance, and local name.

The meet compatibility row is read before the minimum row because $m$ is meaningful only after the common surface $M$ has synchronized the two source schedules. The sibling packet \autoref{ch:concrete-instances-cauchymodulusmeet-namecert} supplies exactly this shared-modulus boundary.

The finite minimum exactness row is then a dyadic comparison inside $D$. It chooses no infinite stream representative and invokes no ambient completeness theorem; it records only the displayed smaller threshold for each finite request.

Finally $W,Q,E,H,C,P,N$ replay the same route for stream, regular-rational, and Real-facing consumers. Since none of the carrier coordinates stores quotient stream equality, selected limits, host real equality, countable choice, or ambient completeness, the listed rows exhaust the local certificate surface.
\end{proof}

\begin{theorem}[Regular Cauchy minimum-modulus L10 route]\label{thm:regular-cauchy-min-modulus-l10-route}
Let $R=(S_0,S_1,M,m,W,Q,D,E,H,C,P,N)$ be an accepted $\RegularCauchyMinModulusUp$ carrier. Every L10 Real-and-Completion consumer of the minimum schedule factors through the finite route
$$
S_0,S_1 \longmapsto M \longmapsto m,D \longmapsto W,Q \longmapsto E \longmapsto H,C,P,N .
$$
The route first reads the two regular-Cauchy schedules, then the shared $\CauchyModulusMeetUp$ surface, then the finite minimum row with its $\DyadicRatCoreUp$ tolerance ledger, then the $\StreamNameUp$ windows and $\RegSeqRatUp$ readback, and only then reaches the terminal $\RealUp$ seal. It exports no quotient stream equality, selected limit, hidden completion object, countable-choice selector, or ambient $\RealUp$ completeness principle.
\end{theorem}
\begin{proof}
Start with the accepted carrier and the obligation theorem above. They force every consumer request to retain $S_0,S_1,M,m,W,Q,D,E,H,C,P,N$ as one finite packet.

The shared-modulus step reads $M$ before any minimum schedule is used. Thus the consumer is not comparing unrelated streams; it is reading the two source schedules through the same displayed meet surface and the same transported tolerance boundary.

The minimum step reads $m$ together with $D$. The row $D$ carries the dyadic comparisons and tolerance budgets, while $m$ records only the finite threshold selected by those displayed comparisons.

The stream and regular-rational handoff is then mediated by $W$ and $Q$, and the terminal seal by $E$. The structural rows $H,C,P,N$ transport, replay, source, and name the same route. Hence every L10 consumer factors through the displayed sequence, with no quotient stream, selected limit, hidden completion object, countable choice, or ambient completeness coordinate available.
\end{proof}

\falsifiablePrediction{An accepted $\RegularCauchyMinModulusUp$ consumer that reaches a $\RegSeqRatUp$ readback or $\RealUp$ seal without reading $S_0,S_1,M,m,W,Q,D,E,H,C,P,N$ refutes the carrier surface.}

\independenceWitness{cauchymodulusmeet, streamname, regseqrat, dyadicratcore, real: no listed sibling supplies the joint packet containing two regular-Cauchy schedules, a shared meet row, a finite minimum row, stream windows, regular rational readback, dyadic tolerance ledger, terminal Real seal, transport, replay, provenance, and local NameCert row.}

\closureat{RegularCauchyMinModulusUp}{seedStr}

\begin{closurestatus}{\RegularCauchyMinModulusUp}
  \constructivestory{A $\RegularCauchyMinModulusUp$ packet is generated from two regular-Cauchy schedules, a CauchyModulusMeet row, a finite minimum row, StreamName windows, RegSeqRat readback, a DyadicRatCore tolerance ledger, a Real seal, componentwise hsame transport, displayed Cont replay, Pkg provenance, and a local NameCert row.}
  \theoryclosure{\seedClosure}
  \scopeclosed{\autoref{def:regular-cauchy-min-modulus-carrier}, \autoref{thm:regular-cauchy-min-modulus-namecert-obligations}, and \autoref{thm:regular-cauchy-min-modulus-l10-route} seed the finite minimum-modulus route from two regular-Cauchy schedules through a shared meet surface to stream, regular-rational, and Real-facing consumers.}
  \formalstatus{\formalTargetV}
  \bridgestatus{none}
  \notclaimed{This chapter does not close Lean theorem verification, bridge checking, quotient stream equality, selected limits, hidden completion objects, countable choice, or ambient $\RealUp$ completeness.}
  \upgradepath{Obligation closure requires checked carrier admission, meet compatibility, finite minimum exactness, readback stability, and the L10 Real-and-Completion route for BEDC.Derived.RegularCauchyMinModulusUp.}
\end{closurestatus}
